\section{The selected inner return and its stable sheet}\label{sec:inner-return}
\subsection{The inner section and its stable sheet}

Fix the phase-zero voltage section through the inner periodic orbit
\(\Gamma_i\).  The exact orbit crosses this section with speed bounded below
by \(0.2067539137\).  After removing the translation multiplier, the
section return has one simple real unstable multiplier and a stable
complement.  The validated characteristic root satisfies
\begin{equation}
 0.6983604100<s_u<0.6983604300,                         \label{eq:inner-unstable-root}
\end{equation}
while the stable spectral radius is at most \(e^{-0.01}\).

Let \(q\) be the unstable right column and let \(\widetilde f\) be the raw
left functional supplied by the bordered problem.  Once
\(\widetilde f(q)\ne0\), put
\(f=\widetilde f/\widetilde f(q)\).  Write the section coordinates as
\[
 u=f(\phi),\qquad s=\phi-qf(\phi).
\]
A forward Lyapunov--Perron construction seeks the local stable sheet in the
form
\begin{equation}
 u=\psi(s),\qquad \psi(0)=D\psi(0)=0.                  \label{eq:stable-graph}
\end{equation}
The graph transform is estimated in six directed blocks corresponding to
stable and unstable inputs and outputs.  Its difficult term is the stable
part of the second return variation in two unstable directions,
\[
 Y_{qq}-qf(Y_{qq}).
\]
Bounding \(Y_{qq}\) and \(qf(Y_{qq})\) separately loses the cancellation and
does not close the graph radius.

\subsection{The continuous RFDE adjoint functional}

For a discrete-delay variational equation
\[
 \dot z(t)=A_0(t)z(t)+\sum_j A_j(t)z(t-\tau_j),
\]
the periodic advanced adjoint carries a current atom and an absolutely
continuous history density.  Its bilinear history pairing, not a Hermitian
inner product, defines the raw section functional \(\widetilde f\).  The Fourier
cokernel-to-history reconstruction has the exact reversal
\begin{equation}
 \widehat r_n=E_{-,-n},                                \label{eq:fourier-reversal}
\end{equation}
with no complex conjugation.  This distinction is essential for the
annihilation of the neutral tangent.

The directed tail solve proves
\begin{equation}
 \text{tail contraction}<0.104433459418,\qquad
 \norm{r_{\rm tail}}_{\ell^1}<0.01150531427,\qquad
 \norm r_{\ell^1}<0.1835103672.                        \label{eq:adjoint-tail}
\end{equation}
The Grushin border derivative and the phase-averaged RFDE bilinear identity
then give the directed normalization
\begin{equation}
 0.000305331752871<\abs{\widetilde f(q)}
 <0.000546927481134.                                   \label{eq:fq-bound}
\end{equation}
Thus the Fourier cokernel, continuous periodic adjoint, and Route--C
history functional can be normalized to the same object \(f\).  A recovery-only
history shard has normalized action greater than \(12.9099\), which also
rules out any current-voltage-only projection.

For a \(240\)-step finite-section diagnostic, performing the correlated
deflation before taking the norm gives
\begin{equation}
 \norm{Y_{qq}-qf_N(Y_{qq})}_\infty=7.26112<12,          \label{eq:yqq-pilot}
\end{equation}
whereas the separate triangle estimate is \(56.7830\).  This is strong
evidence for the required correlation, but \cref{eq:yqq-pilot} itself is
not interval evidence.

\begin{proposition}[Base-orbit correlated \(uu\) block]\label{prop:base-uu}
Let \(q^\Sigma\) be the physical Route--C section unstable direction and
let \(Y_{qq}\) be the physical-time second returned variation at the
validated center inner orbit.  Then
\begin{align}
 \norm{V_{qq}-\widehat V_{qq}}_P
   &\le0.00808459995104,\notag\\
 \abs{\widetilde f(q^\Sigma)}
   &\ge4.2563852678\times10^{-4},\notag\\
 \norm{Y_{qq}-q^\Sigma
       \widetilde f(Y_{qq})/\widetilde f(q^\Sigma)}_Y
   &\le0.0475498854764,\notag\\
 \norm{q^\Sigma}_Y&\ge0.0775543158981.                 \label{eq:directed-base-uu}
\end{align}
Consequently the normalized stable-output block satisfies
\begin{equation}
 C^{uu}_{s,\mathrm{base}}
 \le\frac{0.0475498854764}{(0.0775543158981)^2}
 <7.905649079<12.                                     \label{eq:base-uu-bound}
\end{equation}
\end{proposition}

The proof uses \(1042\) physical method-of-steps cells with both delays
aligned and a degree-\(24\), \(192\)-bit Taylor--Bernstein residual.  It
defines the quotient through the same-row identity
\[
 \frac{\widetilde f(Y_{qq})}{\widetilde f(q^\Sigma)}-c
 =\frac{\widetilde f(Y_{qq}-cq^\Sigma)}
        {\widetilde f(q^\Sigma)}
\]
and applies the adjoint error only after the center history has been
deflated.  A direct atom--density action on the polynomial guide, the
advanced covariance \(e^sT\), continuous voltage/recovery source residuals,
all cell seams, the density-product variation, and the normalization defect
give a total correlated action error at most
\(1.500297667\times10^{-6}\).  In particular, no exact inhomogeneous
adjoint identity is silently assigned to the approximate Fourier row.

\Cref{prop:base-uu} closes the former covector/cancellation bottleneck only
at one periodic base orbit.  Its uniform split-ball inflation and the other
five blocks---hence all six uniform split-ball Hessian blocks---the split
nonlinear return tube, and the first-positive-return/no-earlier-hit enclosure
remain unvalidated.  The discrete stable-power constant is supplied separately
by \cref{prop:inner-terminal-stable-row}.  Accordingly, the six-block
quantitative stable graph \cref{eq:stable-graph} remains open; a qualitative
local stable graph for the selected coordinate map is proved below in
\cref{prop:split-bimeasure-local-graph}.

The exact inflation budget shows what a uniform proof must achieve.  Since
the base value leaves margin \(4.0943509217\), it is sufficient to prove a
split-ball Lipschitz constant strictly below
\[
 L_{uu}^{\rm req}=2408.4417186722.
\]
A source-bound \(192\)-bit scalar \(P\)-logarithmic-norm attempt closes the
self-consistent enclosure on all \(1042\) physical cells, but first loses
the radius-\(0.01\) local section ball on cell \(581\), where its voltage
radius is already \(0.010059402\); its period-end bound is \(2.151313\).
Thus this positive scalar majorant is rigorously rejected as the uniform
inflation route.  It does not rule out a correlated bound.  Indeed,
\[
 2\tau_0<T<\tau_0+\tau_1<3\tau_0
\]
reduces the signed linear ingress to only the four Volterra words
\(\varnothing,(\tau_0),(\tau_1),(\tau_0,\tau_0)\).  The required replacement
is to form \(U(t)P_s=U(t)-U(t)qf/f(q)\) across those words before taking
total variation.

\begin{proposition}[Primitive ingress and physical-frame obstruction]
\label{prop:inner-primitive-ingress}
On the \(1042\)-cell delay-aligned physical grid, the degree-\(24\),
\(192\)-bit Taylor--Bernstein residuals of the five primitives
\(F,G,C_0,C_1,C_{00}\) are bounded above respectively by
\[
 7.30427\times10^{-7},\quad
 2.54960\times10^{-7},\quad
 3.33270\times10^{-7},\quad
 2.56054\times10^{-7},\quad
 3.06759\times10^{-11}.
\]
After every Hermite seam and coefficient tail is included, the moving-frame
maximum error radii are at most
\[
 4.92742\times10^{-4},\quad
 2.42414\times10^{-4},\quad
 1.73778\times10^{-6},\quad
 1.43306\times10^{-6},\quad
 3.64819\times10^{-11}.
\]
The resulting directed primitive ingress into the event-row estimate is at
most
\begin{equation}
 1.463546184\times10^{-3}.                             \label{eq:primitive-ingress}
\end{equation}
By contrast, the declared unprojected scalar physical-frame Gronwall
propagation returns \(F\)- and \(G\)-error radii larger than
\(3.04997\times10^7\) and \(1.89994\times10^7\), respectively, and
therefore cannot certify the stable row.
\end{proposition}

\begin{proposition}[Selected terminal stable row]
\label{prop:inner-terminal-stable-row}
Let
\[
 Y=C([-\tau_{\max},0],\mathbb R)\times\mathbb R,
 \qquad \Sigma_0=\{h\in Y:h_v(0)=0\},
\]
and set
\[
 A=\Pi_TU(T,0)|_{\Sigma_0},\qquad
 f=\widetilde f/\widetilde f(q^\Sigma),\qquad
 P_s=\Id-q^\Sigma f,\qquad E_s=\ker f.
\]
Here \(\Pi_Ty=y-\dot X_Ty_v(0)/\dot v(T)\) is the terminal section
projection.
For the selected near-one-period phase-fixed linear map,
\begin{equation}
 \norm{AP_s}_{\Sigma_0\to Y}
 \le0.009896427481610001<0.1.                 \label{eq:inner-terminal-stable-row}
\end{equation}
Moreover,
\[
 AP_s=P_sA=P_sAP_s,
 \qquad AP_s(\Sigma_0)\subset E_s.
\]
Consequently, for \(A_s=A|_{E_s}\),
\begin{equation}
 \norm{A_s^n}_{E_s\to Y}\le0.1^n,
 \qquad K_s=1\quad(n\ge1).                   \label{eq:inner-stable-power}
\end{equation}
\end{proposition}

The certificate forms the common signed row
\(R_\theta\Pi_TU(T,0)(\Id-q^\Sigma f)\) before every modulus.  The exact
section quotient removes the current-voltage \(\delta_0\) atom, while the
current-recovery atom remains active.  The true-period enclosure gives
\(T>\tau_{\max}\), so the complete returned history lies after time zero and
has no unadvanced identity block.  Its only method-of-steps words are
\(\varnothing,(\tau_0),(\tau_1),(\tau_0,\tau_0)\).  The continuous calculation
covers \(641\) returned-history phase rows plus one recovery-output row against
\(640\) input-history cells, hence \((641+1)\times640=410{,}880\) support
rectangles.  The independently outward-rounded ledger is
\begin{align*}
 E_{\rm ctr}&\le0.004362999710080374,
 &E_{\rm bin}&\le0.000010000000000001,\\
 E_{\rm 4I}&\le0.001463546183857267,
 &E_{\rm rank}&\le0.004059226965383500,\\
 E_{\rm ratio}&\le6.506172913144082\times10^{-7},
 &E_T&\le4.004997545987388\times10^{-9}.
\end{align*}
The displayed terms sum to at most \(0.009896427481610003\); the unrounded
directed calculation gives the sharper bound in
\cref{eq:inner-terminal-stable-row}.  The certified primitive tubes enter only
through their direct algebraic image: neither the unknown terminal norm nor
\(K_s\) is used in its own estimate.

\Cref{prop:inner-primitive-ingress} is therefore only the primitive parent
budget.
\Cref{prop:inner-terminal-stable-row} supplies the missing common terminal
row, true-\(T_+\) coverage, section quotient, and continuous output-phase
supremum, and thereby closes the discrete one-step stable-map norm.

A stronger intermediate-time stable-flow certificate could instead be
obtained without assuming the stable-flow constant that it is meant to prove.
Let \(U_s(t,s)=U(t,s)P_s(s)\), and let
\(\widehat U_s\) be one common signed polynomial guide whose initial defect
and complete method-of-steps residual lie in the exact transported stable
bundle.  The latter includes the differential residual, history-transport
boundary, delay activations, and cell seams.  The complete-history norm also
retains the unadvanced translation/identity block; it is not only a norm of
the propagated density.  For \(\alpha\ge0\), put
\begin{align*}
 K&=\sup_{s\le t}e^{\alpha(t-s)}\norm{U_s(t,s)},&
 \widehat K&=\sup_{s\le t}e^{\alpha(t-s)}
                    \norm{\widehat U_s(t,s)},\\
 \Delta&=\sup_s\norm{D_s}
 +\sup_{s\le t}\int_s^t e^{\alpha(r-s)}
                    \norm{R(r,s)}\,dr .
\end{align*}
Duhamel's formula remains in the stable bundle and yields
\begin{equation}
 K\le\widehat K+K\Delta,\qquad
 \Delta<1\Longrightarrow
 K\le K_{\rm int}:=\frac{\widehat K}{1-\Delta}.       \label{eq:projected-radii}
\end{equation}
Let \(\widehat M_s\) be the corresponding terminal event-corrected row, let
\(C_{\Pi,T}\) bound the exact event correction, let \(\Delta_T\) be the
preprojection endpoint residual budget, and let \(\varepsilon_{\Pi,T}\)
contain the remaining event-row and speed-denominator errors.  Then
\begin{equation}
 \norm{M_s}\le\norm{\widehat M_s}
                 +C_{\Pi,T}K_{\rm int}\Delta_T
                 +\varepsilon_{\Pi,T}.
                                                               \label{eq:stable-terminal-radii}
\end{equation}
The event projection cannot be assigned norm one; its sampled factor is
about \(2.01358\).  A strict right side below \(\rho_s<1\) would provide the
corresponding intermediate-time stable-flow estimate.  The exact projection
is essential: propagating the five primitives separately violates the
stable-residual premise and recovers the rejected unstable growth.  The phase
correction is applied only at the terminal section; a moving \(\Pi_t\) would
add a derivative term to the residual equation.  This route is not used by
\cref{prop:inner-terminal-stable-row}, whose fixed-start terminal calculation
already proves the discrete power estimate.  At present the common full
\((s,t)\) projected residual, \(\Delta\), and
\(C_{\Pi,T},\Delta_T,\varepsilon_{\Pi,T}\) remain unvalidated.  Thus
\cref{eq:projected-radii,eq:stable-terminal-radii} are an optional route toward
intermediate-flow, nonlinear-tube, or no-earlier-hit control, not a
prerequisite for \(K_s=1\).

\subsection{Eventually smooth event Hessian and two-return graph budget}

The discrete estimate in \cref{prop:inner-terminal-stable-row} does not by
itself make the moving complete-history return \(C^2\) on a ball of arbitrary
continuous histories.  We first record the exact event-aligned Hessian, so
that the remaining analytic and numerical gates are explicit.

Let \(S_t(x)\) be the fixed-time complete-history solution map, let \(T(x)\)
be a selected crossing of the phase section, and put
\[
 a(x)=\dot v_x(T(x)),\qquad \ell_0(y)=y_v(0).
\]
For initial directions \(h,k\), let \(U_h\) and \(V_{hk}\) be the first and
second fixed-time variations, with \(V_{hk,0}=0\), and use the superscript
\(T\) for the complete segment translated to the selected event.  Set
\begin{align}
 n_h&=\ell_0(U_h^T),&
 T_h&=-\frac{n_h}{a},                                      \label{eq:event-hessian-first}\\
 Z_{hk}&=V_{hk}^T-
 \frac{\dot U_h^Tn_k+\dot U_k^Tn_h}{a}
 +\frac{\ddot X_Tn_hn_k}{a^2},&
 T_{hk}&=-\frac{\ell_0(Z_{hk})}{a}.                       \label{eq:event-hessian-second}
\end{align}
Define the moving event projection
\begin{equation}
 \Pi_x y=y-\dot X_T\frac{\ell_0(y)}{a(x)}.                 \label{eq:moving-event-projection}
\end{equation}

\begin{proposition}[Selected-return Hessian]
\label{prop:selected-return-hessian}
On any common \(C^2\) selected-event tube on which \(a(x)>0\), the moving
complete-history return \(P(x)=S_{T(x)}(x)\) satisfies
\begin{equation}
 DP(x)h=\Pi_xU_h^T,\qquad
 D^2P(x)[h,k]=\Pi_xZ_{hk}.                              \label{eq:selected-return-hessian}
\end{equation}
The phase projection in \cref{eq:selected-return-hessian} is applied exactly
once.  It precedes the fixed spectral deflation \(P_s\) or the fixed
unstable row \(\widehat f\); those two projections are not interchangeable.
\end{proposition}

\begin{proof}
Differentiate the scalar event identity twice to obtain
\cref{eq:event-hessian-first,eq:event-hessian-second}, including every
terminal-history translation term.  Substitution into the second derivative
of \(S_{T(x)}(x)\) gives \cref{eq:selected-return-hessian}.  In particular,
\(\ell_0(DP\,h)=\ell_0(D^2P[h,k])=0\), which also checks that no event
correction has been duplicated.
\end{proof}

For the one-period center return, \(T>\tau_{\max}\) removes the unadvanced
identity block from the \emph{linear} output, but \(T<2\tau_{\max}\).  The
old one-period argument therefore does not satisfy the \(C^2\) complete-
segment condition in \cref{thm:eventual-selected-return}.  The selected
near-two-period center does: the frozen orbit intervals give
\begin{equation}
 T_- -2\tau_{\max,+}>14.010740123253945.                \label{eq:two-return-smoothing-margin}
\end{equation}
This proves the center-orbit margin.  The next proposition closes the common
full-ball event inequality on an explicit, deliberately microscopic scale.

\begin{proposition}[Explicit common near-two-period event tube]
\label{prop:explicit-two-period-event-tube}
Let
\[
 B_\lambda=\{(x_s,x_u):\norm{x_s}_Y\le0.0097\lambda,
                    \ \abs{x_u}\le0.00025\lambda\},
 \qquad \lambda_0=9\times10^{-31}.
\]
There is an open reduced-history neighborhood \(W_{\rm open}\subset Y\) and
an open coordinate restriction \(D_{\rm open}\) containing
\(B_{\lambda_0}\) on which every injected history, including every
continuous reduced stable history
satisfying the displayed norm bound, has one unique positively oriented
voltage event in the common physical-time window
\begin{equation}
 I_2=[36.3714198982518419857,\;36.3734198982598419858].
                                                               \label{eq:explicit-two-period-window}
\end{equation}
The event time and the complete reduced-history hit are \(C^2\); the initial
history and the hit both lie in the fixed radius-\(0.01\) local phase-zero
section patch.  Uniformly on the larger open comparison scale
\(9.1\times10^{-31}\), the endpoint gap is greater than
\(1.03751\times10^{-4}\), the event speed is greater than
\(0.244880\), the returned-history radius is less than
\(5.09376\times10^{-4}\), and
\(T_--2\tau_{\max}>14.010740\).

Writing \(\chi\) for the fixed section chart, the induced map has type
\(P_{\rm sel}:D_{\rm open}\to D_{\rm out}\), where
\(D_{\rm out}=\chi(\Sigma_{\rm loc})\).  No containment
\(D_{\rm out}\subset D_{\rm open}\) is asserted.

At the center, \(T(Y_*)=2P\) and the hit is exactly \(Y_*\).  The result is
neither a self-map of \(B_{\lambda_0}\), nor a certificate at \(\lambda=1\),
nor an event-ordinal statement.
\end{proposition}

\begin{proof}
Write \(B=\sup_t\abs{v_*(t)-1}\).  On an orbit-centered tube of radius
\(\beta\), the corrected fast-field row gives
\[
 H_\beta=2(1+B+\beta)+12\varepsilon\kappa_3(B+\beta),\qquad
 L_\beta=\max\{A_0+H_\beta\beta,2\varepsilon\}.
\]
Choose \(\beta\) to be half the certified center endpoint gap and propagate
the single complete-history majorant
\[
 M(t)\le\rho_{\rm open}+\int_0^tL_\beta M(s)\,ds
 \qquad(0\le t\le T_+).
\]
It retains the arbitrary initial-history translates before each delay
activates and covers all later method-of-steps seams.  Directed evaluation
gives
\(0.00995(9.1\times10^{-31})e^{L_\beta T_+}<\beta\); the corresponding
open-scale ceiling is greater than \(9.13306649\times10^{-31}\).  A first-exit
argument therefore proves the common flow tube.  The center endpoint signs
and the Lipschitz perturbation bound give the stated strict signs and speed,
so the selected zero is unique in \(I_2\).  The complete-history estimate
also gives terminal section-patch containment.  Finally, apply
\cref{thm:eventual-selected-return} first with the affine compatible lift
\(\mathcal I:W_{\rm open}\to X\), the global polynomial RFDE field, and
\(g_X=g_Y\circ\pi\).  Projection gives the ambient reduced \(C^2\) hit, and
restriction through the section injection \(j:D_{\rm open}\to
W_{\rm open}\cap\Sigma_{\rm loc}\) gives the stated coordinate map.  The
strict margin in \cref{eq:two-return-smoothing-margin} supplies the required
joint smoothing.
\end{proof}

The preceding scalar comparison is correct but extremely cancellation blind.
The recovery equation supplies a weighted energy in which the two
instantaneous voltage--recovery cross terms cancel exactly.  Exploiting that
identity enlarges the same selected-event domain by more than twenty decimal
orders.

\begin{proposition}[Orbit-aware enlarged near-two-period event tube]
\label{prop:orbit-aware-two-period-event-tube}
With the same anisotropic balls \(B_\lambda\) and the same physical window
\cref{eq:explicit-two-period-window}, every history in
\begin{equation}
 B_{1.2\times10^{-8}}
\end{equation}
has one unique positively oriented selected voltage event.  The conclusion
holds on the strictly larger open scale \(1.25\times10^{-8}\); the event time
and complete reduced-history hit are \(C^2\), and the hit remains in the
declared radius-\(0.01\) terminal section patch.  Uniformly on that open
scale, the endpoint gap is greater than \(1.05563\times10^{-4}\) and the
event speed is greater than \(0.206347\).

The sufficient construction has scale ceiling
\begin{equation}
 \lambda_{\rm en}>1.2770076307683837\times10^{-8}.
                                                               \label{eq:orbit-aware-lambda-ceiling}
\end{equation}
It proves neither the preferred scale \(\lambda=1\), a same-ball self-map,
an event ordinal, any one of the six Hessian caps, nor a pulse crossing.
\end{proposition}

\begin{proof}
Let \(x=\eta_v(t)\) and \(y=\eta_w(t)\) be the difference from the exact
periodic orbit and set
\[
 z(t)^2=x(t)^2+\varepsilon^{-1}y(t)^2.
\]
If \(a(t)\) is the current variational coefficient and \(b_0(t),b_1(t)\)
are the two delayed coefficients, then
\begin{equation}
 \frac12\frac{d}{dt}z^2
 =a(t)x^2-y^2+\sum_{j=0}^1b_j(t)x(t)x(t-\tau_j)+x(t)N(t).  \label{eq:weighted-energy-cancellation}
\end{equation}
The terms \(-xy+yx\) have cancelled exactly.  For the complete-history
running envelope \(Z\), including the untranslated initial voltage history
before each delay activates, this gives at every increasing-envelope time
\[
 D^+Z(t)\le\left[r(t)+\frac12H_\beta Z(t)\right]Z(t),
 \qquad
 r(t)=\max\{0,\max(a(t),-\varepsilon)+|b_0(t)|+|b_1(t)|\}.
\]
Here \(H_\beta\) is the corrected polynomial Hessian row in the proof of
\cref{prop:explicit-two-period-event-tube}.

On the \(1042\) delay-aligned cells, \(192\)-bit Taylor--Bernstein
arithmetic gives
\begin{equation}
 \int_0^{P_{\rm bin}}r(t)\,dt
 \le6.353397548471557,
 \qquad \sup r(t)\le0.774979654297297.                \label{eq:orbit-aware-rate}
\end{equation}
The enclosure uses the exact binary interval for \(P_{\rm bin}\), rescales
by the directed true-period interval, and adds the phase drift between the
exact delays \(4\sqrt5,5\sqrt5\) and their binary centers before bounding
the delayed coefficients.  Thus neither decimal round-tripping nor a mesh
maximum is used as exact data.

For a reduced-\(Y\) initial radius \(\rho\),
\(Z(0)\le\sqrt{1+\varepsilon^{-1}}\rho=\sqrt6\rho\).
Choose \(\beta\) as half the center endpoint gap.  On a hypothetical first
exit through \(Z=\beta\), integration of the preceding Dini inequality over
two true periods and the remaining window gives a complete-history gain at
most \(332147.926240866\).  At the open scale \(1.25\times10^{-8}\), the
resulting deviation is at most
\[
1.011906217152045\times10^{-4}<
1.033769568740104\times10^{-4}=\beta,
\]
a contradiction.  The polynomial
field is bounded on this closed tube, so continuation through \(T_+\) also
holds.  The endpoint signs, speed and terminal patch follow from the same
deviation bound, and \cref{thm:eventual-selected-return} supplies the stated
\(C^2\) event and complete-history hit.
\end{proof}

\begin{proposition}[Direct two-period derivative]
\label{prop:direct-two-period-derivative}
Let \(Q_Y\) be the direct selected hit map supplied by
\cref{prop:explicit-two-period-event-tube} on the reduced physical section,
and let
\[
 A=\Pi U_Y(P,0)|_{\Sigma_0}
\]
be the proved one-period phase-fixed derivative.  Then
\begin{equation}
 Q_Y(Y_*)=Y_*,\qquad DQ_Y(Y_*)=A^2.                    \label{eq:direct-two-period-derivative}
\end{equation}
This conclusion does not require a nonlinear one-period return or an
identity \(Q=P^2\).  If \(J(x_s,x_u)=x_s+\widehat qx_u\) is the section
coordinate isomorphism, then
\begin{equation}
 DQ_{\rm coord}(0)=J^{-1}A^2J.                         \label{eq:two-period-coordinate-conjugacy}
\end{equation}
For the compatible full-history lift \(\mathcal I\), the correctly typed
identity is
\begin{equation}
 DQ_X(\mathcal I(Y_*))D\mathcal I=D\mathcal I\,A^2,   \label{eq:two-period-full-intertwining}
\end{equation}
not \(DQ_X=A^2\).
\end{proposition}

\begin{proof}
Put \(M=U_Y(P,0)\) and
\(\Pi=I-v\ell/a\), where \(v\) is the phase tangent,
\(\ell=Dg\), and \(a=\ell(v)\).  Periodicity and the forward variational
cocycle give \(U_Y(2P,0)=M^2\) and \(Mv=v\).  Direct differentiation of the
selected event gives \(DQ_Y(Y_*)=\Pi M^2|_{\Sigma_0}\).  Since
\((I-\Pi)y=v\ell(y)/a\) and \(\Pi Mv=0\),
\[
 (\Pi M|_{\Sigma_0})^2=\Pi M\Pi M|_{\Sigma_0}
 =\Pi M^2|_{\Sigma_0}.
\]
The chart chain rule and the exact lift/projection factorization give
\cref{eq:two-period-coordinate-conjugacy,eq:two-period-full-intertwining}.
\end{proof}

The two-return route is also quantitatively favorable.  Use the fixed unit
pair
\[
 \widehat q=q/\norm{q}_Y,\qquad
 \widehat f=\norm{q}_Y f,\qquad P_s=\Id-\widehat q\widehat f,
\]
the split radii \(R_s=0.0097\), \(R_u=0.00025\), and the desired graph radius
\(r=0.0094\).  For output \(i\in\{s,u\}\), let
\(C_{i,ss},C_{i,su},C_{i,uu}\) bound the six projected blocks of
\(D^2\mathcal Q\), with the event projection formed before \(P_s\) or
\(\widehat f\).

\begin{proposition}[Fixed two-step linear splitting]
\label{prop:inner-two-step-splitting}
Let \(A\in\mathcal L(\Sigma_0)\), \(q\in\Sigma_0\), and
\(f\in\Sigma_0^*\) satisfy
\[
 Aq=\mu_uq,\qquad fA=\mu_uf,\qquad f(q)=1,
\]
and put \(P_u=qf\), \(P_s=\Id-qf\), \(E_u=\operatorname{span}\{q\}\),
and \(E_s=\ker f\).  Then these are complementary projections,
\(\Sigma_0=E_s\oplus E_u\), and both spaces are invariant under \(A\).
For the linear two-step operator \(B=A^2\),
\begin{equation}
 BP_s=P_sB=P_sBP_s=(AP_s)^2,
 \qquad BP_u=P_uB=\mu_u^2P_u.                         \label{eq:two-step-intertwining}
\end{equation}
If
\(\|A_s^n\|\le K_s\rho_{s,1}^n\)
and
\(\|(A_u)^{-n}\|\le K_u\rho_{u,1}^n\), then
\begin{equation}
 \|B_s^n\|\le K_s(\rho_{s,1}^2)^n,
 \qquad
 \|(B_u)^{-n}\|\le K_u(\rho_{u,1}^2)^n.              \label{eq:two-step-rates}
\end{equation}
For the source-bound inner-orbit operator of
\cref{prop:inner-terminal-stable-row}, this gives
\begin{equation}
 K_{s,2}=1,\quad \rho_{s,2}\le0.01,
 \qquad
 K_{u,2}=1,\quad
 \rho_{u,2}\le
 0.302183501335053468766049321269.                    \label{eq:model-two-step-rates}
\end{equation}
The last number is an inverse/backward unstable rate; the forward multiplier
has modulus \(\abs{\mu_u}^2>1\).
\end{proposition}

\begin{proof}
The normalization \(f(q)=1\) proves the complementary-projection identities.
The two eigen-relations give
\(AP_s=P_sA=P_sAP_s\) and \(AP_u=P_uA=\mu_uP_u\).
Squaring gives \cref{eq:two-step-intertwining}, while
\(B_s^n=A_s^{2n}\) and, on the one-dimensional unstable space,
\((B_u)^{-n}=\mu_u^{-2n}\Id\), proving
\cref{eq:two-step-rates}.  The numerical values use the proved
\(K_s=1\), \(\rho_{s,1}\le0.1\) and the source-bound
\(\rho_{u,1}\le0.549712198641301273\ldots\).
\end{proof}

This proposition is linear.  If a nonlinear selected map
\(P:D\to\Sigma_0\) is used to define \(\mathcal Q=P\circ P\), its correct
domain is \(D_2=D\cap P^{-1}(D)\); only there does the chain rule give
\(D\mathcal Q(0)=A^2\).  Conversely, a direct near-two-period selected map
has derivative \(A^2\) by \cref{prop:direct-two-period-derivative}, without
first being represented as \(P^2\).  Neither route turns
linear invariance into nonlinear invariance of \(E_s\) or \(E_u\).

\begin{proposition}[Conditional two-return graph arithmetic]
\label{prop:two-return-graph-arithmetic}
Suppose the direct selected near-two-period map extends from the microscopic
domain in \cref{prop:explicit-two-period-event-tube} to the declared split
ball, is \(C^2\) there, and returns that ball to the graph-transform domain.
Identify its coordinates with \(\Sigma_0\) through \(J\) and use the pullback
norm \(\norm{z}_J=\norm{Jz}_Y\).  Then
\cref{prop:direct-two-period-derivative,prop:inner-two-step-splitting} supply
\(D\mathcal Q(0)=B=A^2\),
\(K_{s,2}=K_{u,2}=1\) and the stable forward and unstable inverse/backward
rates
\[
 \rho_{s,2}\le0.01,\qquad
 \rho_{u,2}\le(0.549712198641302)^2,
\]
Take graph weight \(\beta=0.9999\), and suppose the map has the six uniform
Hessian bounds
\begin{equation}
 (C_{s,ss},C_{s,su},C_{s,uu},C_{u,ss},C_{u,su},C_{u,uu})
 \le(1,10,1000,5,10,1000).                            \label{eq:two-return-wide-box}
\end{equation}
Then the exact split Lyapunov--Perron majorant has Perron bound
less than \(0.193086\), stable and unstable self-map slacks greater than
\(1.96408\times10^{-4}\) and \(1.24109\times10^{-4}\), respectively, and
the unit-\(\widehat q\) graph budgets
\begin{equation}
 \norm{\widehat\psi}\le1.25886\times10^{-4},
 \qquad \norm{D\widehat\psi}\le0.026559.              \label{eq:two-return-graph-budgets}
\end{equation}
Thus the graph-transform arithmetic closes under the displayed premises.
The preferred-ball extension, return-domain containment, and six Hessian
bounds remain premises; the center derivative, splitting, and rates do not.
\end{proposition}

A source-bound but nonrigorous signed finite-section pilot preserves the
moving-event quotient, translates the entire returned history at the same
event time, applies the fixed-center \(\widehat q/\widehat f\) pair, and
only then takes the six tensor norms.  The primary
\(n_{\rm step}=240\) row is
approximately
\[
 (7.45637\!\times\!10^{-7},\,4.46663\!\times\!10^{-3},\,29.353924,
   0.594988,\,0.573991,\,158.531687),
\]
and the disclosed last-two-mesh heuristic envelope is
\begin{equation}
 (2.35230\!\times\!10^{-6},\,0.00872968,\,32.117788,
   0.596928,\,0.577596,\,158.531779).                  \label{eq:two-return-pilot-envelope}
\end{equation}
All six entries lie well inside \cref{eq:two-return-wide-box}.  The adapter
at the current-history endpoint is a one-sided four-node stencil; an earlier
zero-future convention is retained only as a discarded audit.  Neither mesh
convergence nor the heuristic envelope is an interval error bound, so
\cref{eq:two-return-pilot-envelope} is design evidence and not a hypothesis
verification.

The obstruction is not merely that three meshes are too few.

\begin{lemma}[Finite-sampling obstruction on an arbitrary \(C^0\) ball]
\label{lem:finite-sampling-obstruction}
Let \(K\) be a nontrivial compact interval,
\(X=C(K)\) with the supremum norm, and let
\(S_Nh=(h(t_1),\ldots,h(t_N))\) sample any finite node set.  For every
bounded linear reconstruction \(R_N:\R^N\to X\),
\begin{equation}
 \norm{\Id-R_NS_N}_{\mathcal L(X)}\ge1.                \label{eq:sampling-no-go}
\end{equation}
Moreover, nodal data cannot determine all unit bilinear forms on \(X\): for
any off-grid \(t_*\), the form
\(B_*(h,k)=h(t_*)k(t_*)y_*\), with \(\norm{y_*}=1\), is invisible to inputs
whose sampled values vanish.
\end{lemma}

\begin{proof}
Choose \(t_*\notin\{t_1,\ldots,t_N\}\) and a continuous spike \(h\) with
\(\norm h_\infty=1\), \(h(t_j)=0\), and \(h(t_*)=1\).  Then \(S_Nh=0\), so
\((\Id-R_NS_N)h=h\), which proves \cref{eq:sampling-no-go}.  Applying the
same spike in both slots proves the bilinear assertion.
\end{proof}

Consequently, the finite-section grid trend cannot be promoted by mesh refinement
alone to an operator-norm bound on the arbitrary continuous-history ball.
A strict route must instead interpret directed tensor coefficients as
physical evaluation atoms and bound the signed atom--density/bimeasure
residual to the true RFDE Hessian.  If \(A=(a_{oij})\) is such a directed
finite tensor and the output reconstruction has norm one, its atomic lift
satisfies
\begin{equation}
 \norm{B_N}\le\max_o\sum_{i,j}\abs{a_{oij}}.           \label{eq:atomic-bimeasure-lift}
\end{equation}
This makes the row-\(\ell^1\) shape of the pilot appropriate, but supplies no
residual bound.  The frozen continuous splitting gives
\(\norm{\widehat f}\le21.810501\) and
\(\norm{P_s}\le22.810501\); hence an ambient Hessian error can be amplified
by as much as \(\norm{P_s}^3\).  The resulting simultaneous raw-error design
ceiling is \(7.58219\times10^{-5}\).  Direct enclosure of each already
projected signed bimeasure is therefore the preferred rigorous route; all six
actual block slots remain open.

The preceding obstruction leaves intact the continuous dual structure of the
RFDE.  The next result both identifies the exact split norm needed by a
validated residual calculation and records the qualitative stable manifold
that follows without a numerical Hessian cap.

\begin{proposition}[Split quotient, bimeasure residual, and local selected graph]
\label{prop:split-bimeasure-local-graph}
Let \(\ell_0(\phi,w)=\phi(0)\), choose a Hahn--Banach extension
\(\bar f\in Y^*\) of \(\widehat f\in\Sigma_0^*\), and write
\[
 E_s=\ker\ell_0\cap\ker\bar f,
 \qquad E_u=\operatorname{span}\{\widehat q\}.
\]
For every \(\mu\in Y^*\), restriction to the stable input space has the
exact quotient norm
\begin{equation}
 \norm{\mu|_{E_s}}
 =\inf_{\alpha,\beta\in\R}
   \norm{\mu-\alpha\ell_0-\beta\bar f}_{Y^*},
 \qquad
 \sup_{|c|\le1}|\mu(c\widehat q)|=|\mu(\widehat q)|.    \label{eq:stable-quotient-cost}
\end{equation}
Here \(Y^*=\mathcal M([-\tau_{\max},0])\times\R\), with the total-variation
plus recovery-atom norm, and the quotient value is independent of the chosen
extension \(\bar f\).
Denote the two input costs in \cref{eq:stable-quotient-cost} by
\(d_s(\mu)\) and \(d_u(\mu)\), respectively, and put
\[
 \omega_s(y)=\norm{y-\widehat q\widehat f(y)}_Y,
 \qquad \omega_u(y)=|\widehat f(y)|.
\]
Here \(I_s:E_s\hookrightarrow\Sigma_0\) and
\(I_u:\R\to\Sigma_0\), \(I_uc=c\widehat q\), are the input injections,
while \(O_s=P_s\) and \(O_u=\widehat f\) are the output maps.  For each
\((o,a,b)\in\{s,u\}^3\), let
\(B_i(z)=O_oD^2Q_Y(Y_*+Jz)[I_a\,\cdot,I_b\,\cdot]\), with the affine
section charts understood and the two mixed-input orders identified, giving
the six directed blocks indexed by \(i\).
If a signed atom--density bimeasure has an absolutely summable or
Bochner-integrable representation
\(\mathcal B=\sum_r y_r\otimes\mu_r\otimes\nu_r\), with
\(y_r\in\Sigma_0\), define
\[
 \mathcal N_{ab}^{o}(\mathcal B)
 =\inf_{\mathcal B=\sum_r y_r\otimes\mu_r\otimes\nu_r}
   \sum_r\omega_o(y_r)d_a(\mu_r)d_b(\nu_r),
 \qquad a,b,o\in\{s,u\}.
\]
Then
\begin{equation}
 \norm{O_o\mathcal B[I_a\,\cdot,I_b\,\cdot]}_{\rm bil}
 \le \mathcal N_{ab}^{o}(\mathcal B).                 \label{eq:split-projective-bound}
\end{equation}
Consequently, if \(C_i\) is an outward split-projective center bound and
\(\delta_{i,r}\) are outward residuals for the base tube and coefficients,
first and second variations, event quotient and history translation,
continuous \(q/f\) output action, and full-ball inflation/return domain, then
\begin{equation}
 C_i+\sum_r\delta_{i,r}<\operatorname{cap}_i
 \quad(i=1,\ldots,6)
 \quad\Longrightarrow\quad
 \sup_{z\in B_\lambda}\norm{B_i(z)}<\operatorname{cap}_i
 \quad(i=1,\ldots,6).                                 \label{eq:six-block-residual-implication}
\end{equation}

Finally, the selected coordinate map
\(Q_{\rm coord}:D\to E_s\times\R\) supplied by
\cref{prop:explicit-two-period-event-tube} has a qualitative local
\(C^2\) stable graph.  More precisely, there are neighborhoods \(U_s\) of
zero in \(E_s\) and \(U\) of zero in \(E_s\times\R\), and a \(C^2\) map
\(\psi:U_s\to\R\) with \(\psi(0)=D\psi(0)=0\), whose graph consists of
the points whose selected-map iterates remain in \(U\) and converge to zero.
This conclusion has no effective domain radius, graph height, or slope.  It
does not identify the graph with the periodic orbit's stable-set germ and
implies neither a pulse crossing nor biological onset.
\end{proposition}

\begin{proof}
The restriction \(Y^*\to E_s^*\) is the metric quotient by
\(E_s^\perp=\operatorname{span}\{\ell_0,\bar f\}\); Hahn--Banach gives
equality in \cref{eq:stable-quotient-cost}.  Apply these exact input costs
and the fixed output split term by term in a signed projective
representation, then take its infimum, to obtain
\cref{eq:split-projective-bound}.  The triangle inequality gives
\cref{eq:six-block-residual-implication}.

For the last assertion,
\cref{prop:direct-two-period-derivative,prop:inner-two-step-splitting} give
the hyperbolic linearization
\(DQ_{\rm coord}(0)=J^{-1}A^2J\), with stable rate at most \(0.01\) and
unstable inverse rate at most \(0.302184\).  Since \(Q_{\rm coord}\) is
\(C^2\) on an open neighborhood, the derivative of its nonlinear remainder
is arbitrarily small after that neighborhood is shrunk.  The local stable-
manifold theorem for Banach-space maps gives the asserted \(C^2\) graph,
without requiring a self-map of the full anisotropic ball.
\end{proof}

The quotient formula removes phase atoms and multiples of \(\bar f\) before
total variation is taken, while the unstable input is evaluated on
\(\widehat q\) with its sign retained.  Thus
\cref{eq:six-block-residual-implication} is sharper than multiplying an
ambient residual by powers of \(\norm{P_s}\), but it supplies no numerical
ingress by itself.  At present the center values and every listed residual
are unset: zero of the six numerical Hessian caps in
\cref{eq:two-return-wide-box} has been proved.  The local stable graph above
therefore cannot be assigned the preferred radius \(R_s=0.0097\), the pulse
radius \(0.0094\), or either quantitative budget in
\cref{eq:two-return-graph-budgets}.

The wide box has the cancellation-blind pair-sum bound \(2000\), so it
violates the older sufficient scalar target \(K_{\rm ret}<188.912224\) while
still closing the correlated split majorant.  Hence that scalar target is not
a necessary graph-transform input; a source-bound direct return-domain proof
may replace it.

There is one indispensable normalization adapter downstream.  The graph
budgets in \cref{eq:two-return-graph-budgets} use the unit-
\(\widehat q\) coordinate, whereas the pulse endpoint intervals use the
physical right column.  From the proved lower bound
\(\alpha_q=\norm{q_{\rm phys}}_Y\ge
\alpha_{q,-}=0.07755431589814009\), the conditional
physical bounds are
\begin{equation}
 \norm{\psi_{\rm phys}}\le0.001623186095,\qquad
 \norm{D\psi_{\rm phys}}\le0.3424510530.              \label{eq:physical-graph-adapter}
\end{equation}
The height in \cref{eq:physical-graph-adapter} exceeds the older convenient
target \(0.001\); that target was sufficient but not maximal.  Direct use of
the endpoint-specific intervals gives the strict margins below.
Combining these with the proved pulse endpoint and derivative enclosures
would give
\begin{equation}
 \begin{aligned}
 H(J_-)&>0.01805400129,&
 H(J_+)&<-0.01316048384,\\
 H'(I_J)&\subset[-260.79414671,-243.20585329].
 \end{aligned}                                        \label{eq:wide-box-conditional-crossing}
\end{equation}
Accordingly, once the proved microscopic event/return tube is enlarged to the
declared graph ball, its graph-transform return domain and
\cref{eq:two-return-wide-box} are certified in the identical chart, and once
either the direct selected-return hypotheses of
\cref{lem:direct-selected-stable-germ} or an identity \(\mathcal Q=P^2\) on
nested domains is verified, a unique \emph{selected} stable-sheet crossing
follows with large strict margin.  The center derivative, fixed splitting,
intertwinings, and linear rates are already proved.  The preferred-domain
quantitative stable graph, crossing, stable-set identification, two-sided
routing, and biological onset remain unproved.  The qualitative selected-map
stable graph in \cref{prop:split-bimeasure-local-graph} does not close any of
these gates.  A no-earlier-hit theorem is needed only for a first-return or
ordinal label.


